\documentclass[12pt,a4paper]{article}
\usepackage[utf8]{inputenc}
\usepackage[T2A]{fontenc}
\usepackage[english,russian]{babel}
\usepackage{amsmath, amssymb, amsfonts, amsthm}
\usepackage{geometry}
\usepackage{float}
\geometry{top=2cm, bottom=2cm, left=2cm, right=2cm}
\usepackage{fancyhdr}
\usepackage{graphicx}
\usepackage{xcolor}
\usepackage{tcolorbox}
\tcbuselibrary{theorems, skins}
\usepackage{tikz}
\usetikzlibrary{calc, arrows.meta, positioning}

% Определение цветов для разных блоков
\definecolor{defcolor}{RGB}{70, 130, 180}   % синий для определений
\definecolor{thmcolor}{RGB}{180, 70, 70}    % красный для теорем
\definecolor{itcolor}{RGB}{60, 180, 100}    % зеленый для IT-примеров
\definecolor{proofcolor}{RGB}{100, 100, 100} % серый для доказательств

% Настройка колонтитулов
\setlength{\headheight}{33pt}
\addtolength{\topmargin}{-18.5pt}
\pagestyle{fancy}
\fancyhf{}
\fancyhead[L]{
\includegraphics[height=1cm]{logo.jpg}}
\fancyhead[C]{\textbf{РТУ МИРЭА, Кафедра ВМ}}
\fancyhead[R]{\thepage}

% Окружения с цветами (без автонумерации заголовков)
\newtcbtheorem[no counter]{definition}{Определение}{
  enhanced, sharp corners, boxrule=1pt, colback=white, colframe=defcolor,
  fonttitle=\bfseries, coltitle=white, attach boxed title to top
  left={yshift=-2mm},
  boxed title style={sharp corners, colback=defcolor}
}{def}

\newtcbtheorem[no counter]{theorem}{Теорема}{
  enhanced, sharp corners, boxrule=1pt, colback=white, colframe=thmcolor,
  fonttitle=\bfseries, coltitle=white, attach boxed title to top
  left={yshift=-2mm},
  boxed title style={sharp corners, colback=thmcolor}
}{thm}

\newtcbtheorem[no counter]{remarkIT}{Применение в IT}{
  enhanced, sharp corners, boxrule=0.8pt, colback=itcolor!10, colframe=itcolor,
  fonttitle=\bfseries, coltitle=white, attach boxed title to top
  left={yshift=-2mm},
  boxed title style={sharp corners, colback=itcolor}
}{it}

\newenvironment{myproof}{\noindent\textbf{Доказательство.}}{\hfill$\blacksquare$}

% Нумерованные списки с отступами
\usepackage{enumitem}

\title{Лекция 11: Обзор поверхностей второго порядка}
\author{}
\date{}

\begin{document}

\maketitle
\thispagestyle{fancy}

\section*{Введение: Зачем IT-специалисту квадрики?}

В предыдущих лекциях мы изучали линии второго порядка на плоскости
(эллипс, гипербола, парабола). Сегодня мы поднимаемся в трёхмерное
пространство и рассматриваем \textbf{поверхности второго порядка}
(квадрики) — множества точек, координаты которых удовлетворяют
алгебраическому уравнению второй степени относительно $x$, $y$, $z$.

Для студента IT-направления квадрики — это не музейные экспонаты из
курса геометрии, а рабочий инструмент. Любая гладкая поверхность в
графике, физике и оптимизации локально аппроксимируется квадрикой,
а сами квадрики дают компактное аналитическое описание форм.

\begin{remarkIT}{Компьютерная графика и геймдев}{gr}
  \begin{itemize}
    \item \textbf{Примитивы и коллизии:} Сфера (частный случай
      эллипсоида), цилиндр и конус — базовые «коллайдеры» в физических
      движках (Unity PhysX, Unreal). Проверка пересечения луча с
      квадрикой сводится к решению квадратного уравнения — это
      основа алгоритмов трассировки лучей (Ray Tracing).
    \item \textbf{Ограничивающие объёмы:} Эллипсоид часто используют
      как bounding volume для ускорения отсечения невидимых объектов
      (frustum culling): проверить точку относительно эллипсоида
      гораздо дешевле, чем относительно сложного меша.
    \item \textbf{Линейчатые поверхности:} Однополостный гиперболоид
      и гиперболический параболоид (седло) строятся из прямых линий
      (прямолинейные образующие). Это используется в процедурной
      генерации архитектуры и в NURBS-моделировании.
  \end{itemize}
\end{remarkIT}

\begin{remarkIT}{Data Science и оптимизация}{ds}
  Квадратичная форма $\Phi(\mathbf{x}) = \mathbf{x}^T A \mathbf{x}$
  задаёт поверхность уровня, тип которой определяется знаками
  собственных значений матрицы $A$:
  \begin{itemize}
    \item Все собственные значения одного знака $\Rightarrow$
      поверхности уровня — \textbf{эллипсоиды}. Функция потерь
      (loss) имеет единственный минимум, градиентный спуск сходится.
    \item Собственные значения разных знаков $\Rightarrow$
      \textbf{седловая точка} (гиперболический параболоид). Это
      главная трудность обучения нейросетей: оптимизатор «застревает»
      в седле, где градиент мал, но это не минимум.
    \item Соотношение полуосей эллипсоида (число обусловленности)
      определяет скорость сходимости: «вытянутый» эллипсоид —
      медленная сходимость, отсюда методы предобуславливания и Adam.
  \end{itemize}
\end{remarkIT}

\begin{remarkIT}{Инженерия, связь и робототехника}{eng}
  \begin{itemize}
    \item \textbf{Антенны и оптика:} Параболоид вращения фокусирует
      параллельный пучок в одну точку — отсюда спутниковые «тарелки»
      и зеркала телескопов.
    \item \textbf{Архитектура и прочность:} Однополостный гиперболоид
      (Шуховская башня, градирни) обладает высокой жёсткостью при
      минимуме материала за счёт линейчатой структуры.
    \item \textbf{SLAM и навигация:} В задачах локализации
      поверхности уровня погрешности оценки положения робота —
      эллипсоиды доверия (covariance ellipsoids).
  \end{itemize}
\end{remarkIT}

Цель лекции — дать каноническую классификацию поверхностей второго
порядка, научиться по уравнению распознавать тип поверхности и
понимать её геометрию через метод сечений.

\subsection*{Алгебраические поверхности и их порядок}

Прежде чем перечислять конкретные формы, дадим строгое определение
того, что мы классифицируем.

\begin{definition}{Алгебраический многочлен степени $n$}{poly_def}
  Функция $\Phi(x, y, z)$ называется \textbf{алгебраическим
  многочленом степени $n$}, если $\Phi(x, y, z)$ есть сумма конечного
  числа членов вида
  \[
    A \cdot x^p \cdot y^q \cdot z^r,
  \]
  где $p + q + r \le n$, и существует по крайней мере одно слагаемое,
  у которого $p + q + r = n$.
\end{definition}

\begin{definition}{Алгебраическая поверхность $n$-го порядка}{surf_def}
  \textbf{Поверхностью}, определяемой в некоторой декартовой системе
  координат алгебраическим многочленом $n$-ой степени, называется
  \textbf{алгебраической поверхностью $n$-ого порядка}. Это множество
  точек $(x, y, z)$, удовлетворяющих уравнению $\Phi(x, y, z) = 0$.
\end{definition}

\begin{definition}{Примеры поверхностей первого и второго
  порядка}{order_examples}
  \textbf{Первый порядок.} Уравнение
  \begin{equation}
    Ax + By + Cz + D = 0 \tag{11.1}
  \end{equation}
  задаёт алгебраическую поверхность 1-го порядка — это
  \textbf{плоскость} (при условии, что не все $A, B, C$ равны нулю).

  \textbf{Второй порядок (общий вид).}
  \begin{equation}
    \begin{aligned}
      \Phi(x, y, z) = {}& a_{11}x^2 + a_{22}y^2 + a_{33}z^2
      + 2a_{12}xy + 2a_{13}xz + 2a_{23}yz \\
      & {}+ 2a_1 x + 2a_2 y + 2a_3 z + a_0 = 0.
    \end{aligned} \tag{11.2}
  \end{equation}
  Это \textbf{алгебраическая поверхность 2-го порядка}. Старшие члены
  образуют квадратичную форму с матрицей
  \[
    A =
    \begin{pmatrix}
      a_{11} & a_{12} & a_{13} \\
      a_{12} & a_{22} & a_{23} \\
      a_{13} & a_{23} & a_{33}
    \end{pmatrix}.
  \]
\end{definition}

\begin{remarkIT}{Канонизация уравнения = диагонализация матрицы}{canon_it}
  Произвольное уравнение второго порядка приводится к каноническому
  виду в два шага: поворот системы координат убирает смешанные члены
  $xy, xz, yz$ (это \textbf{диагонализация} симметричной матрицы $A$
    через её ортогональные собственные векторы — задача о собственных
  значениях), а параллельный перенос убирает линейные члены (выделение
  полного квадрата). Знаки собственных значений матрицы $A$ полностью
  определяют тип квадрики. Именно поэтому численные библиотеки
  (NumPy \texttt{eigh}, Eigen) кладут в основу классификации
  поверхностей алгоритм нахождения спектра.
\end{remarkIT}

Далее перечислим \textbf{канонические} (приведённые к простейшему
виду) уравнения невырожденных и цилиндрических поверхностей второго
порядка. Для исследования формы каждой поверхности используется
\textbf{метод сечений}: поверхность рассекается плоскостями,
параллельными координатным, и анализируются получающиеся плоские
кривые.

\subsection*{Невырожденные поверхности второго порядка}

\subsubsection*{1. Эллипсоид}

\begin{definition}{Эллипсоид}{ellipsoid_def}
  \textbf{Эллипсоидом} называется поверхность, задаваемая в
  канонической системе координат уравнением
  \begin{equation}
    \frac{x^2}{a^2} + \frac{y^2}{b^2} + \frac{z^2}{c^2} = 1, \tag{11.3}
  \end{equation}
  где $a, b, c > 0$ — \textbf{полуоси} эллипсоида.

  В частном случае $a = b = c = R$ получаем \textbf{сферу}:
  \begin{equation}
    x^2 + y^2 + z^2 = R^2. \tag{11.4}
  \end{equation}

  \textbf{Метод сечений.} Сечение плоскостью $z = h$ (при $|h| < c$)
  даёт эллипс
  \[
    \frac{x^2}{a^2} + \frac{y^2}{b^2} = 1 - \frac{h^2}{c^2} > 0,
  \]
  который стягивается в точку при $|h| \to c$. Поверхность
  ограничена: $|x| \le a$, $|y| \le b$, $|z| \le c$.
\end{definition}

\begin{figure}[H]
  \centering
  \includegraphics[width=0.4\textwidth]{/home/shafed/course-reprocessing/new-tex/assets/11-1.png}
  \par\smallskip
  \textit{Рисунок 11.1}
\end{figure}

\begin{remarkIT}{Эллипсоид доверия и bounding volume}{ell_it}
  В фильтре Калмана и SLAM ковариационная матрица $\Sigma$ оценки
  состояния задаёт эллипсоид $\mathbf{x}^T \Sigma^{-1}\mathbf{x} = 1$,
  внутри которого истинное положение лежит с заданной вероятностью.
  Длины полуосей — корни из собственных значений $\Sigma$. В графике
  эллипсоид служит дешёвым ограничивающим объёмом для теста видимости
  и грубой проверки столкновений.
\end{remarkIT}

\subsubsection*{2. Однополостный гиперболоид}

\begin{definition}{Однополостный гиперболоид}{one_hyperboloid_def}
  \textbf{Однополостным гиперболоидом} называется поверхность
  \begin{equation}
    \frac{x^2}{a^2} + \frac{y^2}{b^2} - \frac{z^2}{c^2} = 1. \tag{11.5}
  \end{equation}

  \textbf{Метод сечений.} Сечение любой горизонтальной плоскостью
  $z = h$ — эллипс
  \[
    \frac{x^2}{a^2} + \frac{y^2}{b^2} = 1 + \frac{h^2}{c^2} \ge 1,
  \]
  наименьший («горловой») при $h = 0$. Сечения вертикальными
  плоскостями $x = 0$ или $y = 0$ — гиперболы. Поверхность связна
  (одна полость) и неограниченно простирается вдоль оси $Oz$.
\end{definition}

\begin{figure}[H]
  \centering
  \includegraphics[width=0.4\textwidth]{/home/shafed/course-reprocessing/new-tex/assets/11-2.png}
  \par\smallskip
  \textit{Рисунок 11.2}
\end{figure}

\subsubsection*{3. Двуполостный гиперболоид}

\begin{definition}{Двуполостный гиперболоид}{two_hyperboloid_def}
  \textbf{Двуполостным гиперболоидом} называется поверхность
  \begin{equation}
    \frac{x^2}{a^2} + \frac{y^2}{b^2} - \frac{z^2}{c^2} = -1. \tag{11.6}
  \end{equation}

  \textbf{Метод сечений.} При $z = h$ имеем
  \[
    \frac{x^2}{a^2} + \frac{y^2}{b^2} = \frac{h^2}{c^2} - 1.
  \]
  Правая часть неотрицательна лишь при $|h| \ge c$, поэтому
  поверхность состоит из \textbf{двух отдельных полостей}
  («чаш») при $z \ge c$ и $z \le -c$, с вершинами в точках
  $(0, 0, \pm c)$. В полосе $|z| < c$ точек поверхности нет.
\end{definition}

\begin{figure}[H]
  \centering
  \includegraphics[width=0.4\textwidth]{/home/shafed/course-reprocessing/new-tex/assets/11-3.png}
  \par\smallskip
  \textit{Рисунок 11.3}
\end{figure}

\subsubsection*{4. Конус (второго порядка)}

\begin{definition}{Конус}{cone_def}
  \textbf{Конусом второго порядка} называется поверхность
  \begin{equation}
    \frac{x^2}{a^2} + \frac{y^2}{b^2} - \frac{z^2}{c^2} = 0. \tag{11.7}
  \end{equation}

  \textbf{Метод сечений.} При $z = h$ получаем эллипс
  $\frac{x^2}{a^2} + \frac{y^2}{b^2} = \frac{h^2}{c^2}$, радиус
  которого линейно растёт с $|h|$; при $h = 0$ он вырождается в
  точку — \textbf{вершину} конуса. Конус является асимптотической
  поверхностью для обоих гиперболоидов и состоит из прямолинейных
  образующих, проходящих через вершину.
\end{definition}

\begin{figure}[H]
  \centering
  \includegraphics[width=0.4\textwidth]{/home/shafed/course-reprocessing/new-tex/assets/11-4.png}
  \par\smallskip
  \textit{Рисунок 11.4}
\end{figure}

\begin{remarkIT}{Конус видимости и трассировка лучей}{cone_it}
  Пирамида (frustum) видимости камеры аппроксимируется конусом;
  тест «точка внутри конуса» используется в прожекторном освещении
  (spotlight) и в алгоритмах culling. В трассировке лучей пересечение
  луча $\mathbf{r}(t) = \mathbf{o} + t\mathbf{d}$ с конусом приводит к
  квадратному уравнению относительно $t$ — два корня дают точки входа
  и выхода.
\end{remarkIT}

\subsubsection*{5. Эллиптический параболоид}

\begin{definition}{Эллиптический параболоид}{ell_paraboloid_def}
  \textbf{Эллиптическим параболоидом} называется поверхность
  \begin{equation}
    \frac{x^2}{p} + \frac{y^2}{q} = 2z, \qquad p > 0,\; q > 0. \tag{11.8}
  \end{equation}

  \textbf{Метод сечений.} Сечение $z = h > 0$ — эллипс
  $\frac{x^2}{p} + \frac{y^2}{q} = 2h$ (при $h < 0$ точек нет, при
  $h = 0$ — вершина в начале координат). Сечения вертикальными
  плоскостями $x = 0$, $y = 0$ — параболы, ветвями вверх. Поверхность
  имеет форму чаши и расположена в полупространстве $z \ge 0$.
\end{definition}

\begin{figure}[H]
  \centering
  \includegraphics[width=0.4\textwidth]{/home/shafed/course-reprocessing/new-tex/assets/11-5.png}
  \par\smallskip
  \textit{Рисунок 11.5}
\end{figure}

\begin{remarkIT}{Параболоид: антенны и выпуклая оптимизация}{par_it}
  Параболоид вращения ($p = q$) фокусирует параллельный пучок в один
  фокус — принцип спутниковых тарелок и зеркальных телескопов. В
  оптимизации эллиптический параболоид — это модель выпуклой функции
  потерь: единственный минимум в вершине, к которому гарантированно
  сходится градиентный спуск. Метод Ньютона аппроксимирует любую
  гладкую функцию вблизи минимума именно таким параболоидом.
\end{remarkIT}

\subsubsection*{6. Гиперболический параболоид}

\begin{definition}{Гиперболический параболоид}{hyp_paraboloid_def}
  \textbf{Гиперболическим параболоидом} («седлом») называется
  поверхность
  \begin{equation}
    \frac{x^2}{p} - \frac{y^2}{q} = 2z, \qquad p > 0,\; q > 0. \tag{11.9}
  \end{equation}

  \textbf{Метод сечений.} Сечение $z = h \ne 0$ — гипербола (при
  $h = 0$ — пара пересекающихся прямых). Сечение плоскостью $y = 0$ —
  парабола ветвями вверх ($x^2/p = 2z$), а сечение $x = 0$ — парабола
  ветвями вниз ($-y^2/q = 2z$). Поэтому поверхность имеет характерную
  форму седла с седловой точкой в начале координат.
\end{definition}

\begin{figure}[H]
  \centering
  \includegraphics[width=0.4\textwidth]{/home/shafed/course-reprocessing/new-tex/assets/11-6.png}
  \par\smallskip
  \textit{Рисунок 11.6}
\end{figure}

\begin{remarkIT}{Седловые точки в обучении нейросетей}{saddle_it}
  В пространствах высокой размерности у функции потерь критических
  точек типа «седло» гораздо больше, чем локальных минимумов. Вблизи
  такой точки поверхность локально совпадает с гиперболическим
  параболоидом: по одним направлениям функция растёт, по другим —
  убывает. Градиент здесь близок к нулю, и наивный градиентный спуск
  замедляется. Современные оптимизаторы (Adam, RMSProp, моменты)
  специально проектируются, чтобы «выскакивать» из седловых областей.
\end{remarkIT}

\subsection*{Цилиндрические поверхности}

Цилиндр получается параллельным переносом плоской кривой
(направляющей) вдоль фиксированного направления (образующей). Если
в уравнении отсутствует одна из координат, поверхность —
цилиндрическая с образующими вдоль соответствующей оси. Ниже
образующие параллельны оси $Oz$.

\subsubsection*{7. Эллиптический цилиндр}

\begin{definition}{Эллиптический цилиндр}{ell_cyl_def}
  \textbf{Эллиптическим цилиндром} называется поверхность
  \begin{equation}
    \frac{x^2}{a^2} + \frac{y^2}{b^2} = 1. \tag{11.10}
  \end{equation}
  В сечении любой плоскостью $z = \text{const}$ — один и тот же эллипс
  с полуосями $a$ и $b$; образующие — прямые, параллельные оси $Oz$.
  При $a = b$ получается круговой цилиндр.
\end{definition}

\begin{figure}[H]
  \centering
  \includegraphics[width=0.4\textwidth]{/home/shafed/course-reprocessing/new-tex/assets/11-7.png}
  \par\smallskip
  \textit{Рисунок 11.7}
\end{figure}

\subsubsection*{8. Гиперболический цилиндр}

\begin{definition}{Гиперболический цилиндр}{hyp_cyl_def}
  \textbf{Гиперболическим цилиндром} называется поверхность
  \begin{equation}
    \frac{x^2}{a^2} - \frac{y^2}{b^2} = 1. \tag{11.11}
  \end{equation}
  В каждом сечении $z = \text{const}$ — одна и та же гипербола с
  асимптотами $y = \pm \frac{b}{a} x$; образующие параллельны оси
  $Oz$. Поверхность состоит из двух «желобов».
\end{definition}

\begin{figure}[H]
  \centering
  \includegraphics[width=0.4\textwidth]{/home/shafed/course-reprocessing/new-tex/assets/11-8.png}
  \par\smallskip
  \textit{Рисунок 11.8}
\end{figure}

\subsubsection*{9. Параболический цилиндр}

\begin{definition}{Параболический цилиндр}{par_cyl_def}
  \textbf{Параболическим цилиндром} называется поверхность
  \begin{equation}
    y^2 = 2px. \tag{11.12}
  \end{equation}
  В каждом сечении $z = \text{const}$ — одна и та же парабола;
  образующие параллельны оси $Oz$.
\end{definition}

\begin{figure}[H]
  \centering
  \includegraphics[width=0.4\textwidth]{/home/shafed/course-reprocessing/new-tex/assets/11-9.png}
  \par\smallskip
  \textit{Рисунок 11.9}
\end{figure}

\begin{remarkIT}{Цилиндры как развёртываемые поверхности}{cyl_it}
  Цилиндрические поверхности \textbf{развёртываются} на плоскость без
  растяжения. Это ключевое свойство в текстурировании (UV-mapping):
  цилиндрическая проекция используется для наложения текстур на
  трубы, стволы, тела вращения. В CAD/CAM развёртки цилиндров нужны
  для раскроя листового металла.
\end{remarkIT}

\subsection*{Прямолинейные образующие}

\begin{definition}{Прямолинейная образующая}{ruling_def}
  \textbf{Прямолинейной образующей} поверхности называется прямая
  линия, целиком лежащая на данной поверхности.
\end{definition}

\begin{theorem}{Наличие прямолинейных образующих}{ruling_thm}
  Среди рассмотренных поверхностей второго порядка поверхности под
  номерами \textbf{1} (эллипсоид), \textbf{3} (двуполостный
  гиперболоид) и \textbf{5} (эллиптический параболоид) прямолинейных
  образующих \textbf{не имеют}; все остальные поверхности
  прямолинейные образующие \textbf{имеют}.
\end{theorem}

\begin{myproof}
  Эллипсоид ограничен ($|x| \le a$, $|y| \le b$, $|z| \le c$), а любая
  прямая неограниченна, поэтому целиком уместиться на эллипсоиде не
  может — образующих нет. Двуполостный гиперболоид и эллиптический
  параболоид целиком лежат в одном полупространстве относительно
  некоторой плоскости (для параболоида $z \ge 0$, для двуполостного
  гиперболоида каждая полость отделена слоем $|z| < c$ без точек);
  прямая же, не параллельная этой плоскости, пересекает оба
  полупространства, а параллельная — не удовлетворяет уравнению, что
  исключает образующие.

  Для цилиндров (7–9) образующие очевидны — это сами прямые,
  параллельные оси $Oz$. Конус (4) состоит из прямых, проходящих через
  вершину. Для \textbf{однополостного гиперболоида} (2) и
  \textbf{гиперболического параболоида} (6) существование образующих
  показывается явной факторизацией. Например, для седла
  $\frac{x^2}{p} - \frac{y^2}{q} = 2z$ перепишем разность квадратов:
  \[
    \left(\frac{x}{\sqrt p} - \frac{y}{\sqrt q}\right)
    \left(\frac{x}{\sqrt p} + \frac{y}{\sqrt q}\right) = 2z.
  \]
  Тогда при любом параметре $\lambda \ne 0$ система прямых
  \[
    \begin{cases}
      \dfrac{x}{\sqrt p} - \dfrac{y}{\sqrt q} = \lambda, \\[2mm]
      \dfrac{x}{\sqrt p} + \dfrac{y}{\sqrt q} = \dfrac{2z}{\lambda}
    \end{cases}
  \]
  целиком лежит на поверхности (перемножение уравнений даёт исходное
  уравнение седла). Аналогичная факторизация для однополостного
  гиперболоида
  $\frac{x^2}{a^2} - \frac{z^2}{c^2} = 1 - \frac{y^2}{b^2}$ даёт два
  семейства прямолинейных образующих.
\end{myproof}

\begin{remarkIT}{Линейчатые поверхности в инженерии и графике}{ruled_it}
  Однополостный гиперболоид и седло — \textbf{дважды линейчатые}
  поверхности: через каждую точку проходят две прямые образующие.
  Это используют в строительстве (Шуховские башни, градирни ТЭЦ): из
  прямых стальных балок собирают изогнутую жёсткую оболочку. В
  компьютерной графике линейчатые и седловые куски — базовые
  «патчи» при моделировании поверхностей сплайнами и в архитектурной
  процедурной генерации, где криволинейная форма набирается из прямых
  элементов.
\end{remarkIT}

\subsection*{Итоговая классификация}

Сводная таблица канонических уравнений:
\begin{center}
  \begin{tabular}{l l}
    \hline
    \textbf{Поверхность} & \textbf{Каноническое уравнение} \\
    \hline
    Эллипсоид & $\dfrac{x^2}{a^2}+\dfrac{y^2}{b^2}+\dfrac{z^2}{c^2}=1$ \\[2mm]
    Однополостный гиперболоид &
    $\dfrac{x^2}{a^2}+\dfrac{y^2}{b^2}-\dfrac{z^2}{c^2}=1$ \\[2mm]
    Двуполостный гиперболоид &
    $\dfrac{x^2}{a^2}+\dfrac{y^2}{b^2}-\dfrac{z^2}{c^2}=-1$ \\[2mm]
    Конус & $\dfrac{x^2}{a^2}+\dfrac{y^2}{b^2}-\dfrac{z^2}{c^2}=0$ \\[2mm]
    Эллиптический параболоид & $\dfrac{x^2}{p}+\dfrac{y^2}{q}=2z$ \\[2mm]
    Гиперболический параболоид & $\dfrac{x^2}{p}-\dfrac{y^2}{q}=2z$ \\[2mm]
    Эллиптический цилиндр & $\dfrac{x^2}{a^2}+\dfrac{y^2}{b^2}=1$ \\[2mm]
    Гиперболический цилиндр & $\dfrac{x^2}{a^2}-\dfrac{y^2}{b^2}=1$ \\[2mm]
    Параболический цилиндр & $y^2=2px$ \\
    \hline
  \end{tabular}
\end{center}

Тип поверхности однозначно определяется знаками собственных значений
матрицы квадратичной формы и наличием линейных членов — это и есть
вычислительная основа автоматического распознавания квадрик в
графических и инженерных пакетах.

\end{document}
